\chapter{Aplicaciones conformes y el teorema de la aplicación de
Riemann}
\label{ch:b3:conformal}

Una biyección \omterm{def:b3:holomorphic:holo}{holomorfa} entre dos
dominios transporta todo el análisis complejo de uno al otro: tales
aplicaciones --- \emph{conformes}, porque conservan los ángulos --- son
los isomorfismos del mundo \omterm{def:b3:holomorphic:holo}{holomorfo}.
Este capítulo las clasifica allí donde la clasificación es posible (el
disco, el plano: la clave es el lema de Schwarz, una desigualdad de una
potencia asombrosa), construye la teoría de \omterm{def:b3:topology:compact}{compacidad} de las familias
\omterm{def:b3:holomorphic:holo}{holomorfas} (Montel) y demuestra el
teorema de existencia más profundo de la materia: \emph{todo} subdominio
propio \omterm{def:b3:conformal:simplyconnected}{simplemente conexo} de
$\C$, por dentada que sea su frontera, es
\omterm{def:b3:conformal:conformal}{conformemente} equivalente al disco
unidad. Cerramos con las
funciones armónicas y el
\omterm{thm:b3:conformal:poisson}{núcleo de Poisson}, resolviendo el
problema de Dirichlet en el disco --- el dividendo analítico de la
geometría conforme. En todo el capítulo, $\mathbb D = D(0,1)$ y
$\mathbb H = \{\operatorname{Im}z >
0\}$.

\section{Aplicaciones conformes; transformaciones de Möbius}

\begin{definition}\label{def:b3:conformal:conformal}
Una \emph{aplicación conforme}\index{aplicación conforme} (o
biholomorfismo) entre abiertos es una
biyección \omterm{def:b3:holomorphic:holo}{holomorfa}; su inversa es
automáticamente \omterm{def:b3:holomorphic:holo}{holomorfa}
(el~\cref{cor:b3:residues:openmapping}). Dos dominios son
\emph{conformemente equivalentes} si existe una tal aplicación;
$\operatorname{Aut}(\Omega)$ denota el grupo de las autoaplicaciones
conformes. Donde $f' \neq 0$ --- en todas partes, si $f$ es inyectiva
(demostración del~\cref{cor:b3:residues:openmapping}) ---, la
diferencial es la multiplicación por $f'(z) \neq 0$: una semejanza, de
modo que las aplicaciones conformes conservan los ángulos entre curvas,
orientación incluida.
\end{definition}

\begin{example}[Transformaciones de Möbius]
\label{ex:b3:conformal:mobius}
Para $\bigl(\begin{smallmatrix}a & b\\ c &
d\end{smallmatrix}\bigr) \in GL_2(\C)$, la \emph{transformación de
Möbius}\index{transformación de Möbius} $z \mapsto
\frac{az + b}{cz + d}$ es conforme de $\C\setminus\{-d/c\}$ sobre
$\C\setminus\{a/c\}$ (la inversa es del mismo tipo, dada por la matriz
inversa; la composición corresponde al producto de matrices). La
\emph{transformación de Cayley}\index{transformación de Cayley}
\[
\varphi(z) = \frac{z - \iu}{z + \iu}
\]
lleva $\mathbb H$ \omterm{def:b3:conformal:conformal}{conformemente}
sobre $\mathbb D$: en efecto, $\abs{z
- \iu} < \abs{z + \iu}$ exactamente cuando $z$ está más cerca de $\iu$
que de $-\iu$, es decir, $\operatorname{Im}z > 0$; la inversa es
$w \mapsto \iu\frac{1 + w}{1 - w}$. Las transformaciones de Möbius
envían la familia de circunferencias y rectas en sí misma
(el~\cref{exo:b3:conformal:1}).
\end{example}

\begin{example}[La aplicación de
Joukowski]\label{ex:b3:conformal:joukowski}
Más allá de Möbius, la \omterm{def:b3:conformal:conformal}{aplicación
conforme} más útil de la matemática aplicada clásica es
\[
J(z) = \frac12\Bigl(z + \frac1z\Bigr) .
\]
En el exterior $\Omega = \{\abs z > 1\}$ del disco unidad, $J$ es
inyectiva: $J(z) = J(w)$ da $(z - w)(1 -
\frac1{zw}) = 0$ y $\abs{zw} > 1$. Su derivada $J'(z) =
\frac12(1 - z^{-2})$ solo se anula en $z = \pm1$, sobre la frontera: $J$
es una equivalencia conforme de $\Omega$ sobre su imagen, que es
$\C\setminus\intcc{-1}1$ --- la propia circunferencia unidad se pliega
dos a uno sobre el segmento ($J(\eu^{\iu\theta}) = \cos\theta$). Así, el
exterior de un \emph{segmento}, un plano hendido sin frontera regular,
es \omterm{def:b3:conformal:conformal}{conformemente} el exterior de un
disco: las esquinas no son obstáculo para la equivalencia conforme, solo
para la regularidad de la frontera. Las imágenes de circunferencias que
pasan por $\pm1$ pero descentradas son
curvas con forma de perfil alar, y componer $J$ con transformaciones de
Möbius transportó el flujo alrededor de un cilindro --- calculable a
mano --- al flujo alrededor de un ala: durante la primera mitad del
siglo XX, este ejemplo \emph{era} la aerodinámica. Es también la puerta
hacia Chebyshev: $J$ conjuga $z \mapsto z^n$ con el polinomio de
Chebyshev $T_n$ (el~\cref{pb:b3:hilbert:1}, Parte V), pues
$J(z^n) = \cos(n\theta)$ cuando $z = \eu^{\iu\theta}$.
\end{example}

\section{El lema de Schwarz y los grupos de automorfismos}

\begin{theorem}[Lema de Schwarz]\label{thm:b3:conformal:schwarz}
Sea $f \colon \mathbb D \to \mathbb D$
\omterm{def:b3:holomorphic:holo}{holomorfa} con $f(0) = 0$.
Entonces\index{lema de Schwarz}
\[
\abs{f(z)} \leq \abs z \ \ (z \in \mathbb D)
\qquad\text{y}\qquad \abs{f'(0)} \leq 1 ,
\]
y si $\abs{f(z_0)} = \abs{z_0}$ para algún $z_0 \neq 0$, o bien
$\abs{f'(0)} = 1$, entonces $f(z) = \eu^{\iu\theta}z$ es una rotación.
\end{theorem}

\begin{proof}
$g(z) = f(z)/z$ se extiende
\omterm{def:b3:holomorphic:holo}{holomorfamente} a $\mathbb D$ (la
singularidad en $0$ es evitable: $g$ está acotada cerca de $0$,
el~\cref{thm:b3:residues:riemanncw}; su valor en $0$ es $f'(0)$). En
$\abs z = r < 1$: $\abs g \leq \frac1r$, de modo que, por el principio
del máximo (el~\cref{thm:b3:holomorphic:maximum}), $\abs g \leq \frac1r$
en $\bar D(0,r)$; hágase $r \to 1$: $\abs
g \leq 1$ en $\mathbb D$, que son ambas desigualdades. La igualdad en un
punto \omterm{def:b3:topology:interior}{interior} hace que $\abs g$
alcance un máximo \omterm{def:b3:topology:interior}{interior}: $g$
constante de módulo $1$.
\end{proof}

\begin{theorem}[Automorfismos del disco]
\label{thm:b3:conformal:autdisc}
Para $a \in \mathbb D$, el \emph{factor de Blaschke}\index{factor de
Blaschke}
\[
\varphi_a(z) = \frac{z - a}{1 - \bar a z}
\]
es un automorfismo de $\mathbb D$ que intercambia $a$ y $0$, con
$\varphi_a^{-1} = \varphi_{-a}$. Todo automorfismo de $\mathbb D$ es
$\eu^{\iu\theta}\varphi_a$ para $\theta
\in \R/2\pi\Z$, $a \in \mathbb D$ únicos.
\end{theorem}

\begin{proof}
En $\abs z = 1$: $\abs{1 - \bar az} = \abs{\bar z}\abs{1 -
\bar az} = \abs{\bar z - \bar a\abs z^2} = \abs{\bar z - \bar
a} = \abs{z - a}$, de modo que allí $\abs{\varphi_a} = 1$; por el
principio del máximo, $\varphi_a(\mathbb D) \subseteq \bar{\mathbb
D}$, y ser aplicación abierta sitúa la imagen en $\mathbb D$. La
identidad \omterm{def:b3:galois:algebraic}{algebraica}
$\varphi_{-a}\circ\varphi_a = \mathrm{id}$ (cálculo directo) muestra la
biyectividad. Sean ahora $f \in
\operatorname{Aut}(\mathbb D)$ y $a = f^{-1}(0)$: $g =
f\circ\varphi_{-a}$ es un automorfismo que fija $0$. Schwarz aplicado a
$g$ y a $g^{-1}$: $\abs{g(z)} \leq \abs z$ y
$\abs{g^{-1}(w)} \leq \abs w$, de modo que $\abs{g(z)} = \abs z$:
rotación, $g = \eu^{\iu\theta}\,\mathrm{id}$, es decir, $f =
\eu^{\iu\theta}\varphi_a$. Unicidad: $a = f^{-1}(0)$ y $\theta$
evaluando al estilo de $f'$ (o a partir de $f(0) =
-\eu^{\iu\theta}a$ y un valor más).
\end{proof}

\begin{theorem}[Automorfismos del plano]
\label{thm:b3:conformal:autplane}
$\operatorname{Aut}(\C) = \{z \mapsto az + b : a \in \C^*,\ b
\in \C\}$. En consecuencia, $\C$ y $\mathbb D$ no son
\omterm{def:b3:conformal:conformal}{conformemente equivalentes}.
\end{theorem}

\begin{proof}
Sea $f \in \operatorname{Aut}(\C)$ y considérese $g(z) =
f(1/z)$ en $\C^*$: una \omterm{def:b3:holomorphic:holo}{función
holomorfa} con una \omterm{def:b3:residues:singularities}{singularidad
aislada} en $0$. Si fuera esencial, Casorati--Weierstrass
(el~\cref{thm:b3:residues:riemanncw}) haría
$g\bigl(D(0,\varepsilon)\setminus\{0\}\bigr)$ denso, mientras que
$f(D(0, 1))$ es abierto y disjunto de él ($f$ inyectiva: los dos
conjuntos son imágenes de conjuntos disjuntos) --- imposible para un
\omterm{def:b3:topology:interior}{conjunto denso} y un
abierto no vacío. Luego $0$ es un
polo o es evitable para $g$, es decir, $\abs{f(z)}$ tiene a lo sumo
crecimiento polinómico: $f$ es un polinomio
(el~\cref{exo:b3:holomorphic:4}(a)). La inyectividad fuerza el grado
$1$: un polinomio de grado mayor tiene, o bien una raíz múltiple de
$f - c$ en algún sitio ($f'$ se anula), o bien varias preimágenes
distintas (d'Alembert--Gauss, el~\cref{pb:b3:holomorphic:1}); en ambos
casos falla la inyectividad. Por último, una $\C \to
\mathbb D$ conforme sería una función entera acotada: constante
(Liouville) --- no hay equivalencia.
\end{proof}

\section{El teorema de Montel}

\begin{theorem}[Montel]\label{thm:b3:conformal:montel}
Sea $\mathcal F \subseteq \mathcal H(\Omega)$ \emph{localmente acotada}:
todo punto tiene un \omterm{def:b3:topology:topology}{entorno} en el que
$\sup_{f\in
\mathcal F}\sup\abs f < \infty$. Entonces toda sucesión de $\mathcal F$
tiene una subsucesión que converge uniformemente en todos los
\omterm{def:b3:topology:compact}{compactos} de $\Omega$ (con límite
\omterm{def:b3:holomorphic:holo}{holomorfo}).\index{teorema de Montel}
\end{theorem}

\begin{proof}
\omterm{def:b3:complete:equicontinuous}{Equicontinuidad} local: si
$\abs f \leq M$ en $D(a, 2r)
\subseteq \Omega$ para toda $f \in \mathcal F$, la fórmula de Cauchy da,
para $z, z' \in D(a, r)$,
\[
\abs{f(z) - f(z')} = \frac{\abs{z - z'}}{2\pi}
\Bigl|\int_{C_{2r}}\frac{f(w)\,\dd w}{(w-z)(w-z')}\Bigr|
\leq \frac{\abs{z - z'}\,2\pi\cdot2r\,M}{2\pi\,r^2}
= \frac{2M}{r}\,\abs{z - z'} :
\]
una cota de Lipschitz uniforme. Agótese $\Omega$ por
\omterm{def:b3:topology:compact}{compactos} $K_m$; cada $K_m$ queda
recubierto por un número finito de tales discos, de modo que
$\mathcal F$ está uniformemente acotada y es
\omterm{def:b3:complete:equicontinuous}{equicontinua} en $K_m$:
Arzelà--Ascoli (el~\cref{thm:b3:complete:ascoli}) extrae una subsucesión
que converge uniformemente en $K_m$; diagonalícese sobre $m$. El límite
es \omterm{def:b3:holomorphic:holo}{holomorfo} por
\cref{thm:b3:holomorphic:weierstrassconv}.
\end{proof}

\section{El teorema de la aplicación de Riemann}

\begin{definition}\label{def:b3:conformal:simplyconnected}
Un abierto \omterm{def:b3:topology:connected}{conexo}
$\Omega \subseteq \C$ es \emph{simplemente conexo}\index{dominio
simplemente conexo} (en sentido homológico, suficiente para todos
nuestros fines) si $\operatorname{Ind}_\gamma(w) = 0$ para todo ciclo
$\gamma$ en $\Omega$ y todo $w \notin \Omega$ --- «ningún ciclo de
$\Omega$ rodea un agujero». Por el teorema global de Cauchy
(el~\cref{thm:b3:residues:globalcauchy}) y
la~\cref{prop:b3:holomorphic:primitive}, en tal $\Omega$ \emph{toda
\omterm{def:b3:holomorphic:holo}{función holomorfa} tiene primitiva};
por tanto, toda $f \in \mathcal H(\Omega)$ sin ceros tiene un logaritmo
\omterm{def:b3:holomorphic:holo}{holomorfo} ($\exp\circ$(primitiva de
$f'/f$), ajustado por una constante, ya que $(f\eu^{-L})' = 0$) y raíces
$n$-ésimas \omterm{def:b3:holomorphic:holo}{holomorfas}
$\eu^{L/n}$.\index{logaritmo holomorfo}
\end{definition}

\begin{theorem}[Teorema de la aplicación de Riemann]
\label{thm:b3:conformal:rmt}
Todo abierto \omterm{def:b3:conformal:simplyconnected}{simplemente
conexo} $\Omega \subsetneq \C$, $\Omega
\neq \varnothing$, es \omterm{def:b3:conformal:conformal}{conformemente}
equivalente a $\mathbb D$; dado $z_0 \in \Omega$, existe una única
\omterm{def:b3:conformal:conformal}{aplicación conforme} $f \colon
\Omega \to \mathbb D$ con $f(z_0) = 0$ y $f'(z_0) >
0$.\index{teorema de la aplicación de Riemann}
\end{theorem}

\begin{proof}
\emph{Paso 0: la familia no es vacía.} Tómese $b \notin \Omega$: $z - b$
no tiene ceros en $\Omega$, de modo que admite una raíz cuadrada
\omterm{def:b3:holomorphic:holo}{holomorfa} $h$ ($h^2 = z - b$). $h$ es
inyectiva ($h(z) =
h(z')$ elevada al cuadrado da $z = z'$), y si $w \in h(\Omega)$ entonces
$-w
\notin h(\Omega)$ ($h(z) = -h(z')$ también eleva al cuadrado a $z = z'$,
dando $w = -w = 0$, imposible porque $h$ no tiene ceros). Como
$h(\Omega)$ es abierto, contiene un disco $D(h(z_0), \rho)$; entonces
$D(-h(z_0), \rho) \cap h(\Omega) = \varnothing$, es decir,
$\abs{h(z) + h(z_0)} \geq \rho$ para todo $z \in \Omega$. Por tanto,
\[
g(z) = \frac{\rho}{2\,\bigl(h(z) + h(z_0)\bigr)}
\]
es \omterm{def:b3:holomorphic:holo}{holomorfa}, inyectiva (una
\omterm{ex:b3:conformal:mobius}{transformación de Möbius} compuesta con la inyectiva $h$), con
$\abs g \leq \frac12 < 1$. Componiendo con un
\omterm{thm:b3:conformal:autdisc}{factor de Blaschke}
(el~\cref{thm:b3:conformal:autdisc}) para llevar $g(z_0)$ a $0$, la
familia
\[
\mathcal F = \{f \colon \Omega \to \mathbb D \text{
holomorphic, injective, } f(z_0) = 0\}
\]
no es vacía.

\emph{Paso 1: un elemento extremal.} Sea $s =
\sup_{\mathcal F}\abs{f'(z_0)} \in \intoc0{+\infty}$ ($> 0$: sus
miembros son inyectivos, de modo que $f'(z_0) \neq 0$). Tómese $f_n \in
\mathcal F$ con $\abs{f_n'(z_0)} \to s$: la familia está acotada por
$1$, de modo que Montel (el~\cref{thm:b3:conformal:montel}) extrae
$f_n \to f$ uniformemente en los
\omterm{def:b3:topology:compact}{compactos}; $f$ es
\omterm{def:b3:holomorphic:holo}{holomorfa}, $f(z_0) = 0$,
$\abs{f'(z_0)} = s$ (el~\cref{thm:b3:holomorphic:weierstrassconv} para
las derivadas), en particular $f$ no es constante; $f$ es inyectiva por
Hurwitz (el~\cref{exo:b3:residues:8}(b)), y
$f(\Omega) \subseteq \bar{\mathbb D}$, luego $\subseteq
\mathbb D$ (aplicación abierta). Así, $f \in \mathcal F$ alcanza el
supremo: $s < \infty$.

\emph{Paso 2: la aplicación extremal es sobreyectiva.} Supóngase $a \in
\mathbb D\setminus f(\Omega)$. El transporte de Blaschke
$\varphi_a\circ f$ no tiene ceros en el
\omterm{def:b3:conformal:simplyconnected}{simplemente conexo} $\Omega$:
tiene una raíz cuadrada \omterm{def:b3:holomorphic:holo}{holomorfa} $F$
(con $F(\Omega) \subseteq \mathbb D$, pues $\abs F^2 =
\abs{\varphi_a\circ f} < 1$), inyectiva (los cuadrados distinguen).
Normalícese: $G = \varphi_{F(z_0)}\circ F \in
\mathcal F$. Deshaciendo: $f = \varphi_{-a}\circ s_2 \circ
\varphi_{-F(z_0)}\circ G$ donde $s_2(w) = w^2$; la aplicación $\Psi
= \varphi_{-a}\circ s_2\circ\varphi_{-F(z_0)} \colon \mathbb D
\to \mathbb D$ es \omterm{def:b3:holomorphic:holo}{holomorfa} con
$\Psi(0) = f(z_0) = 0$ y \emph{no} es una rotación (no es inyectiva:
$s_2$ no lo es). Lema de Schwarz (caso estricto): $\abs{\Psi'(0)} < 1$,
y la regla de la cadena $f = \Psi\circ G$ da $\abs{f'(z_0)} =
\abs{\Psi'(0)}\,\abs{G'(z_0)} < \abs{G'(z_0)}$ --- en contradicción con
la maximalidad (obsérvese $G \in \mathcal F$). Luego $f$ es
sobreyectiva: una equivalencia conforme.

\emph{Paso 3: normalización y unicidad.} Multiplíquese $f$ por
$\eu^{-\iu\arg f'(z_0)}$ para lograr $f'(z_0) > 0$ (esto no sale de
$\mathcal F$). Si $f_1, f_2$ sirven ambas, $\psi =
f_2\circ f_1^{-1} \in \operatorname{Aut}(\mathbb D)$ fija $0$ con
$\psi'(0) = f_2'(z_0)/f_1'(z_0) > 0$; por
el~\cref{thm:b3:conformal:autdisc}, $\psi$ es una rotación
$\eu^{\iu\theta}$ con $\eu^{\iu\theta} > 0$: $\psi =
\mathrm{id}$.
\end{proof}

\begin{remark}\label{rem:b3:conformal:rmtscope}
El teorema es un enunciado de existencia puro y de alcance asombroso: un
cuadrado, un semiplano, el complementario de un corte, la región entre
dos circunferencias tangentes, un dominio de frontera fractal --- todos
\omterm{def:b3:conformal:conformal}{conformemente} idénticos a
$\mathbb D$. Lo que no da: ninguna fórmula (las aplicaciones explícitas
son la excepción: el~\cref{exo:b3:conformal:5}), ningún comportamiento
en la frontera (de eso se ocupa una teoría más profunda --- el teorema
de Carathéodory), ni unicidad de la extensión a $\C$ ni dominios
múltiplemente \omterm{def:b3:topology:connected}{conexos}: la corona
$\{1 < \abs z < 2\}$ \emph{no} es
\omterm{def:b3:conformal:conformal}{conformemente} un disco punteado, y
coronas con distinta razón de radios son inequivalentes (un hecho
genuinamente más difícil).
\end{remark}

\section{Funciones armónicas y el núcleo de Poisson}

\begin{proposition}\label{prop:b3:conformal:harmonicholo}
Sean $\Omega$ \omterm{def:b3:conformal:simplyconnected}{simplemente
conexo} y $u \colon \Omega \to \R$ armónica\index{función armónica}
($\mathcal C^2$ con $\Delta u = u_{xx} + u_{yy} = 0$). Entonces $u =
\operatorname{Re}F$ para cierta $F$
\omterm{def:b3:holomorphic:holo}{holomorfa}, única salvo constante
imaginaria. En consecuencia, $u$ es $\mathcal C^\infty$, cumple la
propiedad del valor medio y obedece el principio del máximo (ningún
extremo \omterm{def:b3:topology:interior}{interior} estricto salvo si es
constante).
\end{proposition}

\begin{proof}
$g = u_x - \iu u_y$ cumple las
\omterm{prop:b3:holomorphic:cauchyriemann}{ecuaciones de
Cauchy--Riemann} ($P = u_x$, $Q = -u_y$: $P_x = u_{xx} = -u_{yy} = Q_y$
y $P_y = u_{xy} = u_{yx} = -Q_x$) con parciales
\omterm{def:b3:topology:continuity}{continuas}: $g$ es
\omterm{def:b3:holomorphic:holo}{holomorfa}
(la~\cref{prop:b3:holomorphic:cauchyriemann}; la $\R$-diferenciabilidad
se sigue de $\mathcal C^1$). Sea $F_0$ una primitiva (\omterm{def:b3:topology:connected}{conexión} simple,
la~\cref{def:b3:conformal:simplyconnected}); entonces
$\operatorname{Re}F_0$ tiene gradiente $(u_x, u_y)$ ($F_0' = g$ se
desglosa exactamente en eso vía Cauchy--Riemann para $F_0$), de modo que
$u - \operatorname{Re}F_0$ es constante ($\Omega$
\omterm{def:b3:topology:connected}{conexo}): ajústese $F = F_0 + c$. Las
propiedades se transfieren del~\cref{thm:b3:holomorphic:maximum} y
del~\cref{exo:b3:holomorphic:10} (para el principio del máximo aplicado
a $u$ misma, úsese $\eu^{F}$ como allí).
\end{proof}

\begin{theorem}[Fórmula de Poisson; problema de Dirichlet en el
disco]\label{thm:b3:conformal:poisson} Para $0 \leq r < 1$, defínase el
\emph{núcleo de Poisson}\index{núcleo de Poisson}
\[
P_r(\theta) = \sum_{n\in\Z}r^{\abs n}\eu^{\iu n\theta}
= \frac{1 - r^2}{1 - 2r\cos\theta + r^2} \;>\; 0 .
\]
Sea $g \colon \partial\mathbb D \to \R$
\omterm{def:b3:topology:continuity}{continua} y póngase, para
$z = r\eu^{\iu\varphi} \in \mathbb D$,
\[
u(z) = \frac1{2\pi}\int_0^{2\pi}
P_r(\varphi - t)\,g(\eu^{\iu t})\,\dd t .
\]
Entonces $u$ es armónica en $\mathbb D$ y se extiende
\omterm{def:b3:topology:continuity}{continuamente} a $\bar{\mathbb D}$
con valores frontera $g$: es la única
\omterm{prop:b3:conformal:harmonicholo}{función armónica} con esa
propiedad.\index{problema de Dirichlet}
\end{theorem}

\begin{proof}
\emph{Identidades del núcleo}: sumando dos series geométricas,
\[
\sum_{n\in\Z}r^{\abs n}\eu^{\iu n\theta}
= \operatorname{Re}\frac{1 + r\eu^{\iu\theta}}{1 -
r\eu^{\iu\theta}}
= \frac{1 - r^2}{\abs{1 - r\eu^{\iu\theta}}^2},
\]
que es el cociente exhibido; la positividad es clara, y
$\frac1{2\pi}\int_0^{2\pi}P_r = 1$ (solo sobrevive $n = 0$).

\emph{Armonicidad}: con $z = r\eu^{\iu\varphi}$,
\[
u(z) = \operatorname{Re}\biggl[\frac1{2\pi}\int_0^{2\pi}
\frac{\eu^{\iu t} + z}{\eu^{\iu t} - z}\,
g(\eu^{\iu t})\,\dd t\biggr],
\]
(el núcleo entre corchetes tiene parte real $P_r(\varphi - t)$:
calcúlese), y el corchete es \omterm{def:b3:holomorphic:holo}{holomorfo}
en $z$ sobre $\mathbb D$ (el~\cref{exo:b3:holomorphic:7}): $u$ es la
parte real de una \omterm{def:b3:holomorphic:holo}{función holomorfa},
luego armónica.

\emph{Valores frontera}: $P_r(\cdot)$ es una identidad aproximada cuando
$r \to 1^-$: masa $1$, y para $\delta \leq
\abs\theta \leq \pi$, $P_r(\theta) \leq \frac{1 - r^2}{1 -
2r\cos\delta + r^2} \to 0$ uniformemente. La partición habitual
(\omterm{def:b3:topology:continuity}{continuidad} de $g$ cerca de
$\eu^{\iu\varphi_0}$, acotación en el resto) da
$u(r\eu^{\iu\varphi}) \to
g(\eu^{\iu\varphi_0})$ cuando $r\eu^{\iu\varphi} \to
\eu^{\iu\varphi_0}$, uniformemente en el punto frontera: la extensión es
\omterm{def:b3:topology:continuity}{continua}. \emph{Unicidad}: la
diferencia de dos soluciones es armónica en $\mathbb D$,
\omterm{def:b3:topology:continuity}{continua} en la
\omterm{def:b3:topology:interior}{clausura} y nula en la frontera: por
el principio del máximo (aplicado a $\pm$ la diferencia), se anula.
\end{proof}

\begin{omfigure}
\begin{tikzpicture}[scale=1.05]
  \begin{scope}
    \clip (-2.1,-2.1) rectangle (2.1,2.1);
    \foreach \c in {0.2, 0.5, 0.8, 1.1, 1.4}{
      \draw[omDef!60] plot[domain=-2.05:2.05, samples=80]
        (\x, {\c/(2*\x)});
      \draw[omDef!60] plot[domain=-2.05:2.05, samples=80]
        (\x, {-\c/(2*\x)});
      \draw[omThm!60, samples=80, domain=\c:2.9]
        plot ({sqrt((\x*\x-\c*\c)/1)/1}, {0}) ;
    }
    % hyperbolas x^2 - y^2 = c
    \foreach \c in {0.4, 1, 1.6}{
      \draw[omThm!70] plot[domain=-1.32:1.32, samples=60]
        ({cosh(\x)*sqrt(\c)}, {sinh(\x)*sqrt(\c)});
      \draw[omThm!70] plot[domain=-1.32:1.32, samples=60]
        ({-cosh(\x)*sqrt(\c)}, {sinh(\x)*sqrt(\c)});
    }
  \end{scope}
  \draw[->] (-2.2,0) -- (2.3,0);
  \draw[->] (0,-2.2) -- (0,2.3);
  \node[font=\small, omThm] at (1.75,-1.9)
    {$\operatorname{Re}z^2 = c$};
  \node[font=\small, omDef] at (-1.35,1.95)
    {$\operatorname{Im}z^2 = c$};
\end{tikzpicture}

{\small Las curvas de nivel de $\operatorname{Re}z^2 = x^2 - y^2$
(hipérbolas rojas) y de $\operatorname{Im}z^2 = 2xy$ (hipérbolas azules)
se cortan en ángulo recto lejos de $0$: una
\omterm{def:b3:conformal:conformal}{aplicación conforme}
($z \mapsto z^2$, donde $z \neq 0$) conserva la ortogonalidad de las
líneas coordenadas. En $z = 0$, donde la derivada se anula, los ángulos
se duplican.}
\end{omfigure}

\section{Ejercicios}

\begin{exercise}[$\star$]\label{exo:b3:conformal:1}
(a) Verifíquese que la \omterm{ex:b3:conformal:mobius}{transformación de
Cayley} $\varphi(z) = \frac{z -
\iu}{z + \iu}$ es una biyección $\mathbb H \to \mathbb D$ con la inversa
enunciada, y calcúlense las imágenes de $\iu$, $0$, $1$ y $\infty$ (en
el límite). (b) Demuéstrese que $z \mapsto 1/z$ lleva circunferencias y
rectas a circunferencias y rectas. \emph{(Escríbase su ecuación común
$\alpha\abs z^2 + \bar\beta z + \beta\bar z + \gamma = 0$,
$\alpha, \gamma \in \R$.)} Dedúzcase lo mismo para todas las
transformaciones de Möbius.
\end{exercise}

\begin{exercise}[$\star$]\label{exo:b3:conformal:2}
Sea $f \colon \mathbb D \to \mathbb D$
\omterm{def:b3:holomorphic:holo}{holomorfa}. (a) Si $f(0) = 0$ y
$f(a) = a$ para algún $a \neq 0$, demuéstrese que $f
= \mathrm{id}$. (b) Si $f$ es un automorfismo con dos puntos fijos
distintos en $\mathbb D$, demuéstrese que $f = \mathrm{id}$
\emph{(conjúguese por un \omterm{thm:b3:conformal:autdisc}{factor de
Blaschke} para reducirse a (a))}.
\end{exercise}

\begin{exercise}[$\star\star$]\label{exo:b3:conformal:3}
(Schwarz--Pick) Para $f\colon \mathbb D \to
\mathbb D$ \omterm{def:b3:holomorphic:holo}{holomorfa}, demuéstrese
\[
\frac{\abs{f'(z)}}{1 - \abs{f(z)}^2} \;\leq\;
\frac{1}{1 - \abs z^2}
\qquad (z \in \mathbb D),
\]
con igualdad (en un punto, luego en todos) si y solo si $f \in
\operatorname{Aut}(\mathbb D)$. \emph{(Aplíquese Schwarz a
$\varphi_{f(z)}\circ
f\circ\varphi_{-z}$.)} Interpretación: las autoaplicaciones
\omterm{def:b3:holomorphic:holo}{holomorfas} contraen la métrica
hiperbólica.
\end{exercise}

\begin{exercise}[$\star\star$]\label{exo:b3:conformal:4}
Hállense equivalencias conformes explícitas: (a) la banda
$\{0 < \operatorname{Im}z < \pi\} \to \mathbb H$; (b) el cuadrante
$\{\operatorname{Re}z > 0,
\operatorname{Im}z > 0\} \to \mathbb H$; (c) el semidisco
$\mathbb D\cap\mathbb H \to$ a un cuadrante, y después $\to \mathbb H$;
(d) $\mathbb D \to \mathbb D$ que lleve $\frac12$ a $0$ con derivada
positiva allí.
\end{exercise}

\begin{exercise}[$\star\star$]\label{exo:b3:conformal:5}
(a) Demuéstrese que no existe ninguna
\omterm{def:b3:conformal:conformal}{aplicación conforme}
$\C \to \mathbb D$ ni $\C \to
\mathbb H$, ni ninguna $\mathbb D \to \C$. (b) ¿Cuáles de los siguientes
son \omterm{def:b3:conformal:conformal}{conformemente equivalentes} a
$\mathbb D$? Justifíquese mediante el~\cref{thm:b3:conformal:rmt} o
exhibiendo una obstrucción: un cuadrado; $\C\setminus\intoc{-\infty}0$;
$\mathbb D\setminus\{0\}$; $\{1 < \abs z < 2\}$.
\emph{(Para los dos últimos: una imagen conforme del disco punteado se
extendería sobre el punto quitado por
el~\cref{thm:b3:residues:riemanncw}(1) --- desarróllese esto.)}
\end{exercise}

\begin{exercise}[$\star\star$]\label{exo:b3:conformal:6}
Sea $\mathcal F = \{f \in \mathcal H(\mathbb D) : f(0) = 1,\
\operatorname{Re}f > 0\}$. (a) Demuéstrese que $\mathcal F$ está
localmente acotada. \emph{(Compóngase con la aplicación de tipo Cayley
$w \mapsto \frac{w -
1}{w + 1}$ que lleva el semiplano derecho a $\mathbb D$, y aplíquese
Schwarz.)} (b) Dedúzcase la \emph{cota de Herglotz}: $\abs{f(z)} \leq
\frac{1 + \abs z}{1 - \abs z}$ para $f \in \mathcal F$, con sus casos de
igualdad.
\end{exercise}

\begin{exercise}[$\star\star\star$]\label{exo:b3:conformal:7}
¿Dónde usa la demostración del~\cref{thm:b3:conformal:rmt} cada
hipótesis? Sígase la pista de: (i) la \omterm{def:b3:topology:connected}{conexión} simple (dos veces); (ii)
$\Omega \neq \C$; (iii) la \omterm{def:b3:topology:connected}{conexión}.
Demuéstrese después que el teorema \emph{falla} para $\Omega = \C$ y
para la corona, señalando qué paso de la demostración se rompe en cada
caso.
\end{exercise}

\begin{exercise}[$\star\star$]\label{exo:b3:conformal:8}
Resuélvase el problema de Dirichlet en $\mathbb D$ para los datos
frontera: (a) $g(\eu^{\iu t}) = \cos t$; (b) $g(\eu^{\iu t}) =
\cos^2 t$; (c) $g = \mathbf 1_{\text{semicircunferencia superior}}$ --- para (c),
calcúlese $u(0)$ e interprétese vía la propiedad del valor medio.
\emph{(Desarróllese $g$ en serie de Fourier y úsese la serie de $P_r$:
$u(r\eu^{\iu\varphi}) = \sum_n c_n(g)
r^{\abs n}\eu^{\iu n\varphi}$.)}
\end{exercise}

\begin{exercise}[$\star\star\star$]\label{exo:b3:conformal:9}
(Harnack) Sea $u \geq 0$ armónica en $\mathbb D$. Demuéstrese, para
$\abs z = r < 1$:
\[
\frac{1 - r}{1 + r}\,u(0) \;\leq\; u(z) \;\leq\;
\frac{1 + r}{1 - r}\,u(0)
\]
\emph{(acótese el \omterm{thm:b3:conformal:poisson}{núcleo de Poisson}
entre $\frac{1-r}{1+r}$ y $\frac{1+r}{1-r}$; aplíquese la
\omterm{def:b3:representations:rep}{representación} en discos
ligeramente menores y pásese al límite)}. Dedúzcase: una
\omterm{prop:b3:conformal:harmonicholo}{función armónica} en $\C$
acotada inferiormente es constante.
\end{exercise}

\begin{exercise}[$\star\star$]\label{exo:b3:conformal:10}
Usando la invariancia conforme de la armonicidad ($u\circ f$ es armónica
cuando $u$ es armónica y $f$ es
\omterm{def:b3:holomorphic:holo}{holomorfa} --- demuéstrese localmente
vía la~\cref{prop:b3:conformal:harmonicholo}), resuélvase el problema de
Dirichlet en el semiplano superior con datos frontera
$\mathbf 1_{\intoo{-\infty}0}$: demuéstrese que
\[
u(x + \iu y) = \frac1\pi\,\arg(x + \iu y)
\qquad (\arg \in \intoo0\pi \text{ en } \mathbb H)
\]
es armónica en $\mathbb H$ (parte imaginaria de un
\omterm{def:b3:conformal:simplyconnected}{logaritmo holomorfo}) con los
límites frontera requeridos en todo $x \neq
0$, y transpórtese al disco por Cayley para volver a deducir
\cref{exo:b3:conformal:8}(c).
\end{exercise}

\begin{exercise}[$\star\star$]\label{exo:b3:conformal:11}
(Puntos fijos e iteración en el disco) Sea $f\colon\mathbb
D \to \mathbb D$ \omterm{def:b3:holomorphic:holo}{holomorfa}. (a)
Demuéstrese que si $f$ tiene \emph{dos} puntos fijos distintos, entonces
$f = \mathrm{id}$ \emph{(llévese uno a $0$ mediante un automorfismo y
aplíquese el caso de igualdad de Schwarz)}. (b) Supóngase $f(0) = 0$ y
que $f$ no es una rotación. Demuéstrese que las iteradas
$f^{\circ n} \to 0$ uniformemente en todo
\omterm{def:b3:topology:compact}{compacto} $\bar D(0, r)$, $r < 1$
\emph{(Schwarz da $\abs{f(z)}
\leq c_r\abs z$ en $\bar D(0,r)$ con $c_r < 1$ --- justifíquese esta
constante estricta mediante el principio del máximo aplicado a
$f(z)/z$)}.
(c) Ilústrese con $f(z) = \frac{z^2 + z}2$: puntos fijos y velocidad de
convergencia de la \omterm{def:b3:groups:action}{órbita} de $z_0 =
\frac12$.
\end{exercise}

\begin{exercise}[$\star\star$]\label{exo:b3:conformal:12}
(Conjugadas armónicas, en concreto) Sea $u(x, y) = x^3 - 3xy^2
+ 2y$. (a) Compruébese que $u$ es armónica en $\R^2$ y hállense todas
las conjugadas armónicas $v$ (es decir, $u + \iu v$
\omterm{def:b3:holomorphic:holo}{holomorfa}) integrando las
\omterm{prop:b3:holomorphic:cauchyriemann}{ecuaciones de
Cauchy--Riemann}; identifíquese $f(z) =
u + \iu v$ como polinomio en $z$. (b) Demuéstrese que en un
abierto \emph{estrellado} toda
\omterm{prop:b3:conformal:harmonicholo}{función armónica} admite una
conjugada armónica, única salvo constante aditiva \emph{(la $1$-forma
$-u_y\,\dd x +
u_x\,\dd y$ es cerrada; la maquinaria de primitivas
del~\cref{thm:b3:holomorphic:cauchy}, o el lema de Poincaré
del~\cref{ch:b3:forms})}. (c) Dese el contraejemplo estándar en $\C^*$:
$u = \ln\abs z$ no tiene conjugada global --- relaciónese con la forma
angular y el índice
(\cref{ch:b3:forms}).
\end{exercise}

\section{Problema: el teorema del área y el teorema del cuarto
de Koebe}

\begin{problem}[{Problema de fin de semana --- ¿cuánto ha de cubrir
una aplicación univalente?}]\label{pb:b3:conformal:1} Una función
\emph{univalente} es una inyección
\omterm{def:b3:holomorphic:holo}{holomorfa}. Las funciones univalentes
normalizadas en el disco,
\[
\mathcal S = \bigl\{f \in \mathcal H(\mathbb D) \text{
injective},\ f(z) = z + a_2z^2 + a_3z^3 + \cdots\bigr\},
\]
están rígidamente constreñidas: demostraremos la desigualdad de
Bieberbach $\abs{a_2} \leq 2$ y deduciremos el \emph{teorema del cuarto
de Koebe}: la imagen de cualquier $f \in \mathcal S$ contiene el disco
$D(0, \frac14)$ --- la constante universal óptima de la geometría
conforme.\index{teorema del cuarto de Koebe}\index{teorema del área}

\medskip
\noindent\textbf{Parte I --- El teorema del área.} Sea $g(w) = w
+ b_0 + \frac{b_1}w + \frac{b_2}{w^2} + \cdots$
\omterm{def:b3:holomorphic:holo}{holomorfa} e \emph{inyectiva} en
$\{\abs w > 1\}$.
\begin{enumerate}
  \item Para $\rho > 1$, sea $A_\rho$ el área
  (\omterm{def:b3:measure:lebesgueouter}{medida de Lebesgue}) del
  \omterm{def:b3:topology:compact}{compacto} $K_\rho = \C\setminus
        g(\{\abs w > \rho\})$, la región encerrada por la curva de
  Jordan regular $g(C_\rho)$. Usando la fórmula del área del teorema de
  Green--Riemann del volumen de segundo año --- el área encerrada es
  $\frac1{2\iu}
        \oint\bar\zeta\,\dd\zeta$ a lo largo de la frontera orientada
  positivamente ---, demuéstrese que
        \[
        A_\rho = \frac{1}{2\iu}\int_{C_\rho}
        \overline{g(w)}\,g'(w)\,\dd w
        = \pi\Bigl(\rho^2 -
        \sum_{n\geq1}n\,\abs{b_n}^2\rho^{-2n}\Bigr) :
        \]
        sustitúyanse las \omterm{thm:b3:residues:laurent}{series de
        Laurent} de $\bar g$ y $g'$ en $C_\rho$ e intégrese término a
        término (convergencia normal; solo sobreviven los productos de
        frecuencia cero). \item Hágase $\rho \to 1^+$ y conclúyase el
        \emph{teorema del área}:
        \[
        \sum_{n\geq1}n\,\abs{b_n}^2 \;\leq\; 1 .
        \]
        En particular, $\abs{b_1} \leq 1$. ¿Cuándo se da $\abs{b_1}
        = 1$?
\end{enumerate}

\medskip
\noindent\textbf{Parte II --- El $\abs{a_2} \leq
2$ de Bieberbach.} Sea $f = z + a_2z^2 + \cdots \in \mathcal S$.
\begin{enumerate}[resume]
  \item Demuéstrese que $f(z^2)/z^2$ es
  \omterm{def:b3:holomorphic:holo}{holomorfa} y sin ceros en
  $\mathbb D$, y admite una raíz cuadrada
  \omterm{def:b3:holomorphic:holo}{holomorfa} $\varphi$ con
  $\varphi(0) = 1$; póngase $h(z) =
        z\varphi(z^2)$, de modo que $h(z)^2 = f(z^2)$. Demuéstrese que
  $h$ es una función univalente \emph{impar} en $\mathbb
        D$ con desarrollo $h(z) = z + \frac{a_2}2z^3 +
        \cdots$. \emph{(Inyectividad: $h(z)^2 = h(z')^2$ obliga a
  $z^2 = z'^2$; termínese usando la imparidad.)} \item Aplíquese el
  teorema del área a $g(w) = 1/h(1/w) = w -
        \frac{a_2}{2w} + \cdots$ en $\{\abs w > 1\}$ (verifíquense la
  univalencia y el desarrollo) y conclúyase
        $\abs{a_2} \leq 2$.
  \item Demuéstrese que la \emph{función de Koebe}
        \[
        k(z) = \frac{z}{(1 - z)^2} = \sum_{n\geq1}n\,z^n
        \]
        pertenece a $\mathcal S$, tiene $a_2 = 2$ y lleva $\mathbb D$
        sobre $\C\setminus\intoc{-\infty}{-\frac14}$ \emph{(escríbase
        $k = \frac14\bigl[\bigl(\frac{1+z}{1 -
        z}\bigr)^2 - 1\bigr]$ y sígase el rastro de las imágenes)}:
        todas las desigualdades que siguen son óptimas.
\end{enumerate}

\medskip
\noindent\textbf{Parte III --- El teorema del cuarto.}
\begin{enumerate}[resume]
  \item Sean $f \in \mathcal S$ y $c \notin f(\mathbb D)$. Demuéstrese
  que
        \[
        F(z) = \frac{c\,f(z)}{c - f(z)}
        \]
        pertenece a $\mathcal S$, y calcúlese su segundo coeficiente:
        $A_2 = a_2 + \frac1c$. \item Aplíquese Bieberbach a $f$ y a $F$:
        conclúyase
        $\abs{\frac1c} \leq \abs{A_2} + \abs{a_2} \leq 4$,
        es decir, $\abs c \geq \frac14$. \emph{Todo valor omitido tiene
        módulo $\geq \frac14$}: $f(\mathbb D) \supseteq
        D(0, \frac14)$ --- el teorema del cuarto de Koebe. Compruébese
        la optimalidad en la función de Koebe. \item Dedúzcase una
        estimación cuantitativa para la aplicación de Riemann: si
        $\varphi \colon \Omega \to \mathbb D$ es la aplicación de
        Riemann del~\cref{thm:b3:conformal:rmt} en $z_0$, entonces
        \[
        \frac{d\bigl(z_0, \partial\Omega\bigr)}{4}
        \;\leq\; \frac1{\varphi'(z_0)} \;\leq\;
        4\,d\bigl(z_0, \partial\Omega\bigr)
        \]
        --- demuéstrese al menos la desigualdad izquierda aplicando
        Koebe a $\varphi^{-1}$ convenientemente normalizada, y la
        derecha aplicando Schwarz a $\varphi$ en el disco
        $D(z_0, d) \subseteq \Omega$.
\end{enumerate}

\medskip
\noindent\textbf{Parte IV --- Perspectiva.}
\begin{enumerate}[resume]
  \item Bieberbach conjeturó (1916) que $\abs{a_n} \leq n$ para todo
  $n$, con igualdad solo para las rotaciones de la función de Koebe; de
  Branges lo demostró en 1985. Verifíquese la conjetura a mano para la
  función de Koebe y sus rotaciones
  $\eu^{-\iu\theta}k(\eu^{\iu\theta}z)$. Llévese después el desarrollo
  de la pregunta 4 un término más allá: escribiendo
  $h(z) = z + \frac{a_2}2z^3 + c_5z^5 + \cdots$, demuéstrese
  $c_5 = \frac{a_3}2 - \frac{a_2^2}8$ y
        \[
        g(w) = w - \frac{a_2}{2}\,w^{-1} +
        \Bigl(\frac{3a_2^2}8 - \frac{a_3}2\Bigr)w^{-3} +
        \cdots ,
        \]
        de modo que el teorema del área produce la desigualdad afinada
        $\bigl|\frac{a_2}2\bigr|^2 + 3\bigl|\frac{3a_2^2}8 -
        \frac{a_3}2\bigr|^2 \leq 1$. Compruébese en la función de Koebe
        ($a_2 = 2$, $a_3 = 3$).
\end{enumerate}

\medskip
\noindent\textbf{Parte V --- El teorema de distorsión.} La desigualdad
de Bieberbach, transportada por los automorfismos a todo el disco,
controla $f'$ en todas partes. Fíjese $f \in \mathcal
S$.
\begin{enumerate}[resume]
  \item (Transformada de Koebe) Para $z_0 \in \mathbb D$, sea
  $\varphi(z) = \frac{z + z_0}{1 + \bar z_0z}$, un automorfismo del
  disco (el~\cref{thm:b3:conformal:autdisc}) con $\varphi(0) = z_0$.
  Demuéstrese que
        \[
        F(z) = \frac{f(\varphi(z)) - f(z_0)}
        {f'(z_0)\,\bigl(1 - \abs{z_0}^2\bigr)}
        \]
        pertenece a $\mathcal S$ \emph{(la univalencia se hereda;
        calcúlese $\varphi'(0) = 1 - \abs{z_0}^2$ y compruébese la
        normalización; recuérdese $f' \neq 0$ para $f$ inyectiva,
        \cref{def:b3:conformal:conformal})}.
  \item Calcúlese el segundo coeficiente $A_2 = \frac12F''(0)$ de $F$:
        \[
        A_2 = \frac12\Bigl[\bigl(1 - \abs{z_0}^2\bigr)
        \frac{f''(z_0)}{f'(z_0)} - 2\bar z_0\Bigr] ,
        \]
        y dedúzcase de Bieberbach (pregunta 4), para $z =
        r\eu^{\iu\theta}$, la \emph{desigualdad fundamental}:
        \[
        \Bigl|\,z\,\frac{f''(z)}{f'(z)} -
        \frac{2r^2}{1 - r^2}\Bigr|
        \;\leq\; \frac{4r}{1 - r^2} .
        \]
  \item Extráigase la parte real:
        \[
        \frac{2r^2 - 4r}{1 - r^2} \;\leq\;
        \operatorname{Re}\Bigl(z\,\frac{f''(z)}{f'(z)}\Bigr)
        \;\leq\; \frac{2r^2 + 4r}{1 - r^2} .
        \]
  \item Demuéstrese que $\frac{\dd}{\dd t}\log\bigl|
        f'(t\eu^{\iu\theta})\bigr| = \frac1t
        \operatorname{Re}\bigl(z\frac{f''(z)}{f'(z)}\bigr)$ en
  $z = t\eu^{\iu\theta}$ \emph{(para una función $\mathcal C^1$ no nula
  $g$ de variable real, $\frac{\dd}{\dd t}\log\abs g =
        \operatorname{Re}(g'/g)$)}, e intégrense las cotas de la
  pregunta 12 a lo largo del radio para obtener el \emph{teorema de
  distorsión}:
        \[
        \frac{1 - r}{(1 + r)^3} \;\leq\; \abs{f'(z)}
        \;\leq\; \frac{1 + r}{(1 - r)^3},
        \qquad \abs z = r .
        \]
  \item Dedúzcase el \emph{teorema de crecimiento}:
        \[
        \frac{r}{(1 + r)^2} \;\leq\; \abs{f(z)} \;\leq\;
        \frac{r}{(1 - r)^2},
        \qquad \abs z = r
        \]
        \emph{(cota superior: intégrese $f'$ en el segmento $[0, z]$;
        cota inferior: si $\abs{f(z)} < \frac14$, el segmento
        $[0, f(z)]$ está contenido en $f(\mathbb D)$ por la pregunta 7;
        tírese de él hacia atrás por $f^{-1}$ ---
        \omterm{def:b3:holomorphic:holo}{holomorfa} por
        el~\cref{cor:b3:residues:openmapping} --- y acótese
        $\abs{f(z)} = \int_\gamma\abs{f'(\zeta)}
        \,\abs{\dd\zeta} \geq \int_0^r\frac{1 - t}{(1 +
        t)^3}\,\dd t$, usando $\abs{\dd\zeta} \geq
        \dd\abs\zeta$)}. \item Verifíquese que la función de Koebe
        alcanza la igualdad en las cuatro cotas, en $z = r$ para las
        superiores y en $z = -r$ para las inferiores: $k$ es
        simultáneamente el miembro más dilatador de $\mathcal S$ y, en
        el antípoda, el más contractivo.
\end{enumerate}

\medskip
\noindent\textbf{Parte VI --- Rigidez extremal.} En todas las
desigualdades vistas hasta ahora, la igualdad identifica la función de
Koebe salvo rotación. Vamos a demostrarlo y luego a cosechar.
\begin{enumerate}[resume]
  \item Supóngase que $f \in \mathcal S$ tiene $\abs{a_2} = 2$.
  Persíguese la igualdad a través de las preguntas 2--4: el teorema del
  área obliga a $g(w) = w + b_0 + \eu^{\iu\alpha}/w$; la imparidad de
  $h$ hace $g$ impar, de modo que $b_0 = 0$; inviértase para hallar $h$
  y después $f$, y conclúyase que
        \[
        f(z) = \eu^{-\iu\theta}k\bigl(\eu^{\iu\theta}z\bigr)
        \quad\text{con } \eu^{\iu\theta} =
        -\eu^{\iu\alpha} :
        \]
        las rotaciones de la función de Koebe son los únicos miembros de
        $\mathcal S$ con $\abs{a_2} = 2$. \item Demuéstrese que si
        $f \in \mathcal S$ omite un valor $c$ con $\abs c = \frac14$
        exactamente, entonces $f$ es una rotación de la función de
        Koebe, e identifíquese $c =
        -\eu^{-\iu\theta}/4$ \emph{(síganse las igualdades por la cadena
        $4 = \abs{1/c} = \abs{A_2 - a_2}
        \leq \abs{A_2} + \abs{a_2} \leq 4$ de la pregunta 7)}: la
        constante del teorema del cuarto solo se alcanza en la familia
        extremal. \item (Coeficientes a bajo precio) Combínese el
        teorema de crecimiento con las
        \omterm{thm:b3:holomorphic:analytic}{estimaciones de Cauchy}
        (el~\cref{thm:b3:holomorphic:analytic}) en la circunferencia
        $\abs z = 1 - \frac1n$ para demostrar
        \[
        \abs{a_n} \;\leq\; \eu\,n^2 \qquad (n \geq 2) .
        \]
        (De Branges, 1985: $\abs{a_n} \leq n$; el factor $\eu n$ es el
        precio de las herramientas elementales.) \item (Recubrimiento de
        subdiscos) Demuéstrese que para todo $0 < r <
        1$,
        \[
        f\bigl(D(0, r)\bigr) \supseteq
        D\Bigl(0, \frac{r}{(1 + r)^2}\Bigr),
        \]
        óptimo para la función de Koebe, y recupérese el teorema del
        cuarto como $r \to 1^-$. \emph{(Los puntos frontera de la imagen
        abierta $f(D(0,r))$ están en $f(\partial D(0,r))$, luego tienen
        módulo $\geq
        \frac{r}{(1+r)^2}$ por la pregunta 14; un segmento de $0$ a un
        punto omitido de módulo menor tendría que cruzar esa frontera.)}
        \item (Koebe en todo punto) Sea $f$ univalente en $\mathbb D$,
        no necesariamente normalizada, y $z_0
        \in \mathbb D$. Demuéstrese
        \[
        \tfrac14\bigl(1 - \abs{z_0}^2\bigr)\abs{f'(z_0)}
        \;\leq\;
        d\bigl(f(z_0), \partial f(\mathbb D)\bigr)
        \;\leq\;
        \bigl(1 - \abs{z_0}^2\bigr)\abs{f'(z_0)}
        \]
        \emph{(izquierda: teorema del cuarto aplicado a la transformada
        de Koebe de la pregunta 10; derecha: Schwarz
        (el~\cref{thm:b3:conformal:schwarz}) aplicado a
        $\psi^{-1}\circ\hat g$, donde $\hat g(w) =
        f^{-1}\bigl(f(z_0) + dw\bigr)$, siendo $d$ la distancia y $\psi$
        un automorfismo del disco que lleva $0$ a $z_0$)}. ¿Por qué no
        es vacío $\partial f(\mathbb D)$? \item Compruébese la pregunta
        20 en $f = k$ en $z_0 = r \in
        \intoo01$: calcúlese $d\bigl(k(r), \partial k(\mathbb
        D)\bigr) = \frac{(1+r)^2}{4(1-r)^2}$ y verifíquese que la
        desigualdad izquierda es una \emph{igualdad}: la función de
        Koebe satura su propio teorema en todo punto de
        $\intoo01$.
  \item (La moraleja) En diez líneas: ¿qué principio único subyace al
  teorema del área, y cómo brotan de él Bieberbach, el teorema del
  cuarto, la distorsión, el crecimiento y el recubrimiento? Compárese
  con el mundo de Schwarz--Pick del~\cref{thm:b3:conformal:schwarz}: en
  ambos, una sola desigualdad
  \omterm{def:b3:topology:interior}{interior} rigidifica toda la
  geometría, y los extremales son únicos salvo rotación.

\end{enumerate}
\medskip
\noindent\textbf{Parte VII --- \omterm{def:b3:topology:compact}{Compacidad}, inversas y una comprobación
con la realidad.}
\begin{enumerate}[resume]
  \item Demuéstrese que la clase $\mathcal S$ es
  \emph{compacta} para la convergencia
  localmente uniforme: está localmente acotada por el teorema de
  crecimiento, luego es normal (el~\cref{thm:b3:conformal:montel}); y un
  límite localmente uniforme de miembros de $\mathcal S$ vuelve a estar
  en $\mathcal S$ \emph{(las normalizaciones pasan al límite por la
  convergencia de Weierstrass de las derivadas; la inyectividad
  sobrevive por Hurwitz, el~\cref{exo:b3:residues:8}, al no ser
  constante el límite)}. ¿Por qué importa esto para problemas extremales
  como el de Bieberbach? \item Para $f \in \mathcal S$, sea
  $g = f^{-1}$, definida cerca de $0$. Demuéstrese que
  $g(w) = w - a_2w^2 + O(w^3)$, de modo que el segundo coeficiente de la
  inversa obedece la misma cota óptima $\abs{A_2} = \abs{a_2} \leq 2$,
  con igualdad exactamente para las funciones de Koebe rotadas. \item
  Determínese para qué $a \in \C$ el polinomio $f(z)
        = z + az^2$ pertenece a $\mathcal S$: demuéstrese que $f$ es
  inyectivo en $\mathbb D$ si y solo si $\abs a \leq \frac12$
  \emph{(factorícese $f(z_1) - f(z_2)$)}. Conclúyase: para los
  polinomios de grado dos, la cota verdadera de los coeficientes es
  $\frac12$, cuatro veces menor que el $2$ de Bieberbach --- los
  extremales de $\mathcal S$ son objetos genuinamente trascendentes, y
  ningún polinomio se les acerca.
\end{enumerate}
\end{problem}
